\chapter{Actualización}
\label{cap:actualizacion}

En todos los modos: haga un \textbf{respaldo manual} justo antes (capítulo~\ref{cap:respaldos}),
traiga el código nuevo con \texttt{git pull}, reconstruya, vuelva a ejecutar la instalación
(aplica las migraciones nuevas) y verifique (capítulo~\ref{cap:verificacion}).

\begin{aviso}[Las migraciones no se deshacen]
	Una vez aplicadas las migraciones de la versión nueva, volver a la versión anterior del
	programa no devuelve la base a su estado anterior. Para volver atrás hay que restaurar el
	respaldo manual hecho antes de actualizar, y se pierde lo hecho desde entonces. Por eso: haga
	el respaldo, compruebe que terminó bien, anote el commit actual y actualice fuera del horario
	de uso.
\end{aviso}

\section{Docker}

Desde \cod{/opt/digaipe/app/deploy} (o \texttt{C:\bs digaipe\bs app\bs deploy}; en Windows,
sin \texttt{sudo}):

\begin{codigo}
docker compose --profile ops run --rm backup
sudo git log -1 --format=%h
sudo git pull
docker compose build
docker compose up -d
docker compose logs setup
\end{codigo}

\texttt{deploy/.env} no está en Git, así que \texttt{git pull} no lo toca. Compare
\texttt{.env.example} con su \texttt{.env} por si la versión nueva trae variables nuevas.

\section{Servidor físico}

\so{Linux}
\begin{codigo}
sudo systemctl stop digaipe digaipe-respaldo
sudo systemd-run --wait --pipe -p User=digaipe-respaldo -p SupplementaryGroups=digaipe -p EnvironmentFile=/etc/digaipe/backup.env /usr/bin/bash /opt/digaipe/app/deploy/backup/backup.sh
cd /opt/digaipe/app
sudo git log -1 --format=%h
sudo git pull
sudo npm ci
sudo npm run build
sudo /usr/local/sbin/digaipe-setup
sudo systemctl start digaipe digaipe-respaldo
\end{codigo}

\so{Windows}
\begin{codigo}
Stop-Service digaipe
Disable-ScheduledTask -TaskPath '\DIGAIPE\' -TaskName Respaldo
& 'C:\Program Files\Git\bin\bash.exe' -c '. /c/digaipe/bin/entorno.sh && bash /c/digaipe/app/deploy/backup/backup.sh'
Set-Location C:\digaipe\app
git log -1 --format=%h
git pull
npm ci
npm run build
powershell -ExecutionPolicy Bypass -File C:\digaipe\bin\digaipe-setup.ps1
Start-Service digaipe
Enable-ScheduledTask -TaskPath '\DIGAIPE\' -TaskName Respaldo
\end{codigo}

\texttt{git pull} no toca los archivos cargados (\texttt{var/uploads} está excluido de Git) ni la
configuración (fuera de la carpeta). Revise si la versión nueva cambió \texttt{deploy/Caddyfile}
o \texttt{deploy/.env.example}, y traslade los cambios a su \texttt{Caddyfile} y a sus archivos
\texttt{.env}. Si cambió la versión requerida de Node.js o de MySQL, actualícela antes.

\section{Volver a la versión anterior}

Si la versión nueva falla y hay que volver atrás:

\begin{enumerate}
	\item Detenga la aplicación y los respaldos.
	\item Vuelva al commit anotado: \texttt{git checkout <commit>} (con \texttt{sudo} en Linux).
	\item Reconstruya: Docker, \texttt{docker compose build}; servidor físico, \texttt{npm ci} y
	      \texttt{npm run build}.
	\item Restaure el respaldo manual hecho antes de actualizar (capítulo~\ref{cap:respaldos}),
	      incluido el paso de volver a ejecutar la instalación.
\end{enumerate}

Cuando corrijan el problema, vuelva a la rama con \texttt{git checkout <rama>} antes del
siguiente \texttt{git pull}.
